% !TEX root = ../paper.tex
\section{Introduction}
\label{sec:introduction}

Granular-material simulations are central to computational physics and
engineering, with applications ranging from geotechnical hazard prediction to
visual effects in games and films. The \emph{Discrete Element Method}
(DEM)~\cite{cundall1979} models granular media as collections of individual
particles interacting through pairwise contact forces, providing a physically
faithful representation of phenomena such as avalanches, landslides, and
snowfall accumulation. However, DEM simulations are computationally demanding:
at each timestep every particle must detect contacts with its neighbours and
resolve penalty-based forces, yielding an $O(N \cdot k)$ cost per frame (where
$k$ is the average neighbour count) that quickly exceeds the real-time budget of
a single CPU core.

Graphics Processing Units (GPUs), with their thousands of lightweight cores and
high-bandwidth memory, are a natural fit for DEM workloads. Each particle can
be mapped to an independent CUDA thread, and the Structure-of-Arrays (SoA)
memory layout ensures coalesced global-memory transactions. Several works have
demonstrated GPU-accelerated particle
simulations~\cite{green2010,poschel2005}, yet most target high-end datacenter
GPUs. Optimizing for \emph{resource-constrained} Maxwell-class hardware (5
Streaming Multiprocessors, 48\,KB shared memory, 64\,K registers per SM)
introduces additional challenges: every register counts, shared-memory budgets
are tight, and kernel launch overhead is a measurable fraction of frame time.

\subsection{The System}

This paper presents \textbf{Angry Santa}, a real-time DEM snowball-avalanche
simulator implemented in CUDA and benchmarked on two Maxwell-class devices: a
laptop NVIDIA GTX~960M (SM~5.0) and an NVIDIA Jetson~Nano (SM~5.3). The main
components are:

\begin{enumerate}
  \item A \textbf{three-tier particle architecture} (static snowpack, dynamic
        shell, rigid-body core) with a probabilistic sigmoid-based capture model
        that converts snowpack particles into shell particles as the snowball
        rolls.
  \item A \textbf{multi-stream GPU pipeline} with fork/join synchronization
        and an \textbf{adaptive kernel-fusion strategy} that switches between a
        grid-based $O(N \cdot k)$ collision path and a single mega-fused
        $O(N^2)$ kernel depending on the active particle count.
  \item A comprehensive set of \textbf{architecture-aware optimizations}
        --- warp-cooperative shared-memory tiling, lock-free atomic capture,
        zero-spill register allocation, gated event recording, sorted-snowpack
        binary-search range limiting --- each analysed with occupancy and
        bandwidth trade-offs specific to Maxwell.
  \item \textbf{Scalar (SISD) and SIMD\,+\,OpenMP CPU baselines} that share
        the same physics pipeline, enabling direct speedup measurements.
\end{enumerate}

\subsection{Paper Organisation}

Section~\ref{sec:physics} describes the DEM physics model, terrain geometry,
and capture algorithm. Section~\ref{sec:architecture} details the three-tier
data architecture and per-frame data flow. Section~\ref{sec:pipeline} presents
the dual-path GPU pipeline. Section~\ref{sec:optimizations} is the core of the
paper, covering all GPU optimizations with quantitative analysis.
Section~\ref{sec:cpu} summarises the CPU reference implementations.
Section~\ref{sec:results} reports experimental results, and
Section~\ref{sec:conclusion} concludes.
